% !TEX encoding = UTF-8 Unicode
\kap{Parameter}
%Nahezu alle Plugin-/Moduleinstellungen werden in VSTForx durch Parameter getätigt. Diese sind in Form von Drehreglern auf der Bühne platzier- und bedienbar. Außerdem bietet VSTForx die Möglichkeit, einmal platzierte Drehregler miteinander zu verbinden (siehe Kapitel \ref{mousecontext} auf Seite \pageref{mousecontext}).
%Dadurch ist es möglich, mehrere Drehregler gleichzeitig zu steuern. 
Every effect and every module is controlled by parameters. All of these parameters have a knob representation which can be placed and controlled on the main view. With the knobs you are able to bind parameters together, so that every parameter changing will be transmitted to its neighbours.

\ukap{Special Parameter Types}
\paragraph{Free}
%Freie Regler sind nicht an ein Plugin oder Modul gebunden.Es lassen sich beliebig viele Freie Regler auf der Bühne platzieren.
Free parameters are not related to any module or plugin in VSTForx. You can add as many free knobs on the view as you want.
\paragraph{Host}
%Host Regler repräsentieren die VST-Pluginparameter von VSTForx. Sie sind von der Hostanwendung aus bedien- und automatisierbar.
The host parameter appears in your host application as VSTForx parameters. These parameters are the connection between the host application and VSTForx. Use these parameters e.g. for host automation.
\ukap{Operators}
%Zusätzlich können Sie, über das Kontextmenü einer Parameterverbindung, Operatoren auf solch eine Verbindung legen. Operatoren verändern das Verhalten einer Parameter Verbindung. So wird zum Beispiel mit dem ''Inverse'' Operator bewirkt, das sich, das Verhalten des Nachbar-Parameters umkehrt. 
Parameter Connection Operators are modificators for parameter value transmissions. To add an operator move the mouse over a connection and open the context menu. Select one operator from the submenu entry \emph{add operator}. 
Some operators are controlled by parameters too. You can access them using the scene browser.
\\\\
Following operators are available:
\begin{itemize}
\item Inverse
\item Exp/Log
\item Log/Exp
\item Offset
\end{itemize}